% =====================================================================
%  Anatomía de un LLM — Fase I — Módulo 6
%  Teoría de la información
%  Compilar: latexmk -r ../latexmkrc M06-informacion.tex
% =====================================================================
\documentclass[11pt,a4paper,twoside,openany]{book}
\usepackage{../estilo-apuntes}

\newcommand{\ApuntesCurso}{Anatomía de un LLM}
\newcommand{\ApuntesModulo}{Módulo 6}

\usetikzlibrary{shapes.geometric,calc,decorations.pathreplacing}

% Anclas de hiperenlace únicas por tema para los enunciados
\makeatletter
\@for\nombre:=definicion,teorema,proposicion,lema,corolario,ejemplo,observacion\do{%
  \expandafter\xdef\csname theH\nombre\endcsname{\noexpand\thechapter.\noexpand\arabic{definicion}}}
\makeatother

% Notación propia del módulo
\newcommand{\KL}[2]{D(#1\,\|\,#2)}
\newcommand{\X}{\mathcal{X}}
\newcommand{\pref}{\pi_{\mathrm{ref}}}
\DeclareMathOperator{\bpb}{bpb}
\DeclareMathOperator*{\argmax}{arg\,max}
\newcommand{\bin}[1]{\textnormal{\texttt{#1}}}

\hypersetup{pdftitle={Módulo 6. Teoría de la información (resumen)},pdfauthor={Anatomía de un LLM}}

\makeatletter
\newlabel{ej:entropia-tres}{{1.5}{2}{}{ejemplo.1.5}{}}
\newlabel{ej:entropia-tres@cref}{{[ejemplo][5][]1.5}{[1][2][]2}}
\newlabel{ej:binaria}{{1.6}{2}{Entropía binaria}{ejemplo.1.6}{}}
\newlabel{ej:binaria@cref}{{[ejemplo][6][]1.6}{[1][2][]2}}
\newlabel{fig:binaria}{{1.1}{3}{Función de entropía binaria del \cref {ej:binaria}. Una moneda muy sesgada es casi previsible y contiene poca información por lanzamiento}{figure.1.1}{}}
\newlabel{fig:binaria@cref}{{[figure][1][1]1.1}{[1][2][]3}}
\newlabel{ej:dos-bits}{{1.11}{4}{}{ejemplo.1.11}{}}
\newlabel{ej:dos-bits@cref}{{[ejemplo][11][]1.11}{[1][4][]4}}
\newlabel{eje:1-aditiva}{{1.1}{5}{}{ejercicio.1.1}{}}
\newlabel{eje:1-aditiva@cref}{{[ejercicio][1][1]1.1}{[1][5][]5}}
\newlabel{eje:1-archivos}{{1.2}{5}{}{ejercicio.1.2}{}}
\newlabel{eje:1-archivos@cref}{{[ejercicio][2][1]1.2}{[1][5][]5}}
\newlabel{eje:1-cambridge}{{1.3}{5}{}{ejercicio.1.3}{}}
\newlabel{eje:1-cambridge@cref}{{[ejercicio][3][1]1.3}{[1][5][]5}}
\newlabel{eje:1-binaria}{{1.4}{5}{}{ejercicio.1.4}{}}
\newlabel{eje:1-binaria@cref}{{[ejercicio][4][1]1.4}{[1][5][]5}}
\newlabel{eje:1-geometrica}{{1.5}{5}{}{ejercicio.1.5}{}}
\newlabel{eje:1-geometrica@cref}{{[ejercicio][5][1]1.5}{[1][5][]5}}
\newlabel{eje:1-funcion}{{1.6}{5}{}{ejercicio.1.6}{}}
\newlabel{eje:1-funcion@cref}{{[ejercicio][6][1]1.6}{[1][5][]5}}
\newlabel{ej:asimetria}{{2.5}{7}{Asimetría}{ejemplo.2.5}{}}
\newlabel{ej:asimetria@cref}{{[ejemplo][5][]2.5}{[1][7][]7}}
\newlabel{fig:venn}{{2.1}{9}{Relaciones entre entropías e información mutua. Cada círculo representa la entropía de una variable; la intersección es la información compartida y la unión, la entropía conjunta}{figure.2.1}{}}
\newlabel{fig:venn@cref}{{[figure][1][2]2.1}{[1][8][]9}}
\newlabel{ej:mutua-bits}{{2.15}{9}{}{ejemplo.2.15}{}}
\newlabel{ej:mutua-bits@cref}{{[ejemplo][15][]2.15}{[1][9][]9}}
\newlabel{eje:2-vf}{{2.1}{10}{}{ejercicio.2.1}{}}
\newlabel{eje:2-vf@cref}{{[ejercicio][1][2]2.1}{[1][10][]10}}
\newlabel{eje:2-direcciones}{{2.2}{10}{}{ejercicio.2.2}{}}
\newlabel{eje:2-direcciones@cref}{{[ejercicio][2][2]2.2}{[1][10][]10}}
\newlabel{eje:2-kl-moneda}{{2.3}{10}{}{ejercicio.2.3}{}}
\newlabel{eje:2-kl-moneda@cref}{{[ejercicio][3][2]2.3}{[1][10][]10}}
\newlabel{eje:2-conjunta}{{2.4}{10}{}{ejercicio.2.4}{}}
\newlabel{eje:2-conjunta@cref}{{[ejercicio][4][2]2.4}{[1][10][]10}}
\newlabel{eje:2-galletas}{{2.5}{10}{}{ejercicio.2.5}{}}
\newlabel{eje:2-galletas@cref}{{[ejercicio][5][2]2.5}{[1][10][]10}}
\newlabel{eje:2-gibbs-ln}{{2.6}{11}{}{ejercicio.2.6}{}}
\newlabel{eje:2-gibbs-ln@cref}{{[ejercicio][6][2]2.6}{[1][11][]11}}
\newlabel{eje:2-mutua}{{2.7}{11}{}{ejercicio.2.7}{}}
\newlabel{eje:2-mutua@cref}{{[ejercicio][7][2]2.7}{[1][11][]11}}
\newlabel{ej:clases}{{3.3}{12}{}{ejemplo.3.3}{}}
\newlabel{ej:clases@cref}{{[ejemplo][3][]3.3}{[1][12][]12}}
\newlabel{fig:arbol-kraft}{{3.1}{13}{El código prefijo $C_4$ del \cref {ej:clases} como árbol binario: las palabras son las hojas y ninguna está en el camino hacia otra. La suma de Kraft es $\tfrac 12+\tfrac 14+\tfrac 18+\tfrac 18=1$}{figure.3.1}{}}
\newlabel{fig:arbol-kraft@cref}{{[figure][1][3]3.1}{[1][13][]13}}
\newlabel{ej:huffman}{{3.11}{15}{}{ejemplo.3.11}{}}
\newlabel{ej:huffman@cref}{{[ejemplo][11][]3.11}{[1][15][]15}}
\newlabel{fig:huffman}{{3.2}{15}{Árbol de Huffman del \cref {ej:huffman}. Cada nodo interno lleva la suma de las probabilidades de sus hojas}{figure.3.2}{}}
\newlabel{fig:huffman@cref}{{[figure][2][3]3.2}{[1][15][]15}}
\newlabel{eje:3-dos-huffman}{{3.1}{16}{}{ejercicio.3.1}{}}
\newlabel{eje:3-dos-huffman@cref}{{[ejercicio][1][3]3.1}{[1][16][]16}}
\newlabel{eje:3-diadica}{{3.2}{16}{}{ejercicio.3.2}{}}
\newlabel{eje:3-diadica@cref}{{[ejercicio][2][3]3.2}{[1][16][]16}}
\newlabel{eje:3-huffman5}{{3.3}{17}{}{ejercicio.3.3}{}}
\newlabel{eje:3-huffman5@cref}{{[ejercicio][3][3]3.3}{[1][17][]17}}
\newlabel{eje:3-unica}{{3.4}{17}{}{ejercicio.3.4}{}}
\newlabel{eje:3-unica@cref}{{[ejercicio][4][3]3.4}{[1][17][]17}}
\newlabel{eje:3-shannon}{{3.5}{17}{}{ejercicio.3.5}{}}
\newlabel{eje:3-shannon@cref}{{[ejercicio][5][3]3.5}{[1][17][]17}}
\newlabel{eje:3-bloques}{{3.6}{17}{}{ejercicio.3.6}{}}
\newlabel{eje:3-bloques@cref}{{[ejercicio][6][3]3.6}{[1][17][]17}}
\newlabel{eje:3-no-comprime}{{3.7}{17}{}{ejercicio.3.7}{}}
\newlabel{eje:3-no-comprime@cref}{{[ejercicio][7][3]3.7}{[1][17][]17}}
\newlabel{ej:aab}{{4.3}{19}{}{ejemplo.4.3}{}}
\newlabel{ej:aab@cref}{{[ejemplo][3][]4.3}{[1][18][]19}}
\newlabel{fig:aritmetica}{{4.1}{19}{Codificación aritmética del mensaje «aab» (\cref {ej:aab}). Cada nivel subdivide el intervalo del nivel anterior en proporción a las probabilidades del símbolo siguiente}{figure.4.1}{}}
\newlabel{fig:aritmetica@cref}{{[figure][1][4]4.1}{[1][19][]19}}
\newlabel{ej:aab-palabra}{{4.6}{20}{}{ejemplo.4.6}{}}
\newlabel{ej:aab-palabra@cref}{{[ejemplo][6][]4.6}{[1][19][]20}}
\newlabel{fig:compresor}{{4.2}{21}{Compresión con un modelo de lenguaje. Emisor y receptor usan el mismo modelo; el receptor recupera cada token y lo añade al contexto para obtener la distribución del siguiente}{figure.4.2}{}}
\newlabel{fig:compresor@cref}{{[figure][2][4]4.2}{[1][20][]21}}
\newlabel{ej:bpb}{{4.12}{21}{}{ejemplo.4.12}{}}
\newlabel{ej:bpb@cref}{{[ejemplo][12][]4.12}{[1][21][]21}}
\newlabel{eje:4-bpb-concepto}{{4.1}{22}{}{ejercicio.4.1}{}}
\newlabel{eje:4-bpb-concepto@cref}{{[ejercicio][1][4]4.1}{[1][22][]22}}
\newlabel{eje:4-bac}{{4.2}{22}{}{ejercicio.4.2}{}}
\newlabel{eje:4-bac@cref}{{[ejercicio][2][4]4.2}{[1][22][]22}}
\newlabel{eje:4-descodificar}{{4.3}{22}{}{ejercicio.4.3}{}}
\newlabel{eje:4-descodificar@cref}{{[ejercicio][3][4]4.3}{[1][22][]22}}
\newlabel{eje:4-abbb}{{4.4}{22}{}{ejercicio.4.4}{}}
\newlabel{eje:4-abbb@cref}{{[ejercicio][4][4]4.4}{[1][22][]22}}
\newlabel{eje:4-bpb}{{4.5}{22}{}{ejercicio.4.5}{}}
\newlabel{eje:4-bpb@cref}{{[ejercicio][5][4]4.5}{[1][22][]22}}
\newlabel{eje:4-equivocado}{{4.6}{23}{}{ejercicio.4.6}{}}
\newlabel{eje:4-equivocado@cref}{{[ejercicio][6][4]4.6}{[1][22][]23}}
\newlabel{eje:4-equivocado-teo}{{4.7}{23}{}{ejercicio.4.7}{}}
\newlabel{eje:4-equivocado-teo@cref}{{[ejercicio][7][4]4.7}{[1][23][]23}}
\newlabel{ej:monedas}{{5.3}{25}{}{ejemplo.5.3}{}}
\newlabel{ej:monedas@cref}{{[ejemplo][3][]5.3}{[1][25][]25}}
\newlabel{ej:rlhf}{{5.8}{26}{}{ejemplo.5.8}{}}
\newlabel{ej:rlhf@cref}{{[ejemplo][8][]5.8}{[1][26][]26}}
\newlabel{eje:5-beta}{{5.1}{27}{}{ejercicio.5.1}{}}
\newlabel{eje:5-beta@cref}{{[ejercicio][1][5]5.1}{[1][27][]27}}
\newlabel{eje:5-monedas}{{5.2}{27}{}{ejercicio.5.2}{}}
\newlabel{eje:5-monedas@cref}{{[ejercicio][2][5]5.2}{[1][27][]27}}
\newlabel{eje:5-landauer}{{5.3}{28}{}{ejercicio.5.3}{}}
\newlabel{eje:5-landauer@cref}{{[ejercicio][3][5]5.3}{[1][27][]28}}
\newlabel{eje:5-rlhf}{{5.4}{28}{}{ejercicio.5.4}{}}
\newlabel{eje:5-rlhf@cref}{{[ejercicio][4][5]5.4}{[1][28][]28}}
\newlabel{eje:5-lagrange}{{5.5}{28}{}{ejercicio.5.5}{}}
\newlabel{eje:5-lagrange@cref}{{[ejercicio][5][5]5.5}{[1][28][]28}}
\newlabel{eje:5-dpo}{{5.6}{28}{}{ejercicio.5.6}{}}
\newlabel{eje:5-dpo@cref}{{[ejercicio][6][5]5.6}{[1][28][]28}}
\newlabel{ap:soluciones}{{A}{29}{Soluciones de los ejercicios}{appendix.A}{}}
\newlabel{ap:soluciones@cref}{{[appendix][1][2147483647]A}{[1][29][]29}}
\newlabel{sol:eje:1-aditiva}{{A.1}{29}{Tema 1}{section*.4}{}}
\newlabel{sol:eje:1-aditiva@cref}{{[subappendix][1][2147483647,1]A.1}{[1][29][]29}}
\newlabel{sol:eje:1-archivos}{{A.1}{29}{Tema 1}{section*.5}{}}
\newlabel{sol:eje:1-archivos@cref}{{[subappendix][1][2147483647,1]A.1}{[1][29][]29}}
\newlabel{sol:eje:1-cambridge}{{A.1}{30}{Tema 1}{section*.6}{}}
\newlabel{sol:eje:1-cambridge@cref}{{[subappendix][1][2147483647,1]A.1}{[1][30][]30}}
\newlabel{sol:eje:1-binaria}{{A.1}{30}{Tema 1}{section*.7}{}}
\newlabel{sol:eje:1-binaria@cref}{{[subappendix][1][2147483647,1]A.1}{[1][30][]30}}
\newlabel{sol:eje:1-geometrica}{{A.1}{30}{Tema 1}{section*.8}{}}
\newlabel{sol:eje:1-geometrica@cref}{{[subappendix][1][2147483647,1]A.1}{[1][30][]30}}
\newlabel{sol:eje:1-funcion}{{A.1}{31}{Tema 1}{section*.9}{}}
\newlabel{sol:eje:1-funcion@cref}{{[subappendix][1][2147483647,1]A.1}{[1][31][]31}}
\newlabel{sol:eje:2-vf}{{A.2}{31}{Tema 2}{section*.10}{}}
\newlabel{sol:eje:2-vf@cref}{{[subappendix][2][2147483647,1]A.2}{[1][31][]31}}
\newlabel{sol:eje:2-direcciones}{{A.2}{32}{Tema 2}{section*.11}{}}
\newlabel{sol:eje:2-direcciones@cref}{{[subappendix][2][2147483647,1]A.2}{[1][32][]32}}
\newlabel{sol:eje:2-kl-moneda}{{A.2}{32}{Tema 2}{section*.12}{}}
\newlabel{sol:eje:2-kl-moneda@cref}{{[subappendix][2][2147483647,1]A.2}{[1][32][]32}}
\newlabel{sol:eje:2-conjunta}{{A.2}{33}{Tema 2}{section*.13}{}}
\newlabel{sol:eje:2-conjunta@cref}{{[subappendix][2][2147483647,1]A.2}{[1][33][]33}}
\newlabel{sol:eje:2-galletas}{{A.2}{33}{Tema 2}{section*.14}{}}
\newlabel{sol:eje:2-galletas@cref}{{[subappendix][2][2147483647,1]A.2}{[1][33][]33}}
\newlabel{sol:eje:2-gibbs-ln}{{A.2}{34}{Tema 2}{section*.15}{}}
\newlabel{sol:eje:2-gibbs-ln@cref}{{[subappendix][2][2147483647,1]A.2}{[1][34][]34}}
\newlabel{sol:eje:2-mutua}{{A.2}{34}{Tema 2}{section*.16}{}}
\newlabel{sol:eje:2-mutua@cref}{{[subappendix][2][2147483647,1]A.2}{[1][34][]34}}
\newlabel{sol:eje:3-dos-huffman}{{A.3}{35}{Tema 3}{section*.17}{}}
\newlabel{sol:eje:3-dos-huffman@cref}{{[subappendix][3][2147483647,1]A.3}{[1][35][]35}}
\newlabel{sol:eje:3-diadica}{{A.3}{35}{Tema 3}{section*.18}{}}
\newlabel{sol:eje:3-diadica@cref}{{[subappendix][3][2147483647,1]A.3}{[1][35][]35}}
\newlabel{sol:eje:3-huffman5}{{A.3}{36}{Tema 3}{section*.19}{}}
\newlabel{sol:eje:3-huffman5@cref}{{[subappendix][3][2147483647,1]A.3}{[1][36][]36}}
\newlabel{sol:eje:3-unica}{{A.3}{36}{Tema 3}{section*.20}{}}
\newlabel{sol:eje:3-unica@cref}{{[subappendix][3][2147483647,1]A.3}{[1][36][]36}}
\newlabel{sol:eje:3-shannon}{{A.3}{37}{Tema 3}{section*.21}{}}
\newlabel{sol:eje:3-shannon@cref}{{[subappendix][3][2147483647,1]A.3}{[1][36][]37}}
\newlabel{sol:eje:3-bloques}{{A.3}{37}{Tema 3}{section*.22}{}}
\newlabel{sol:eje:3-bloques@cref}{{[subappendix][3][2147483647,1]A.3}{[1][37][]37}}
\newlabel{sol:eje:3-no-comprime}{{A.3}{37}{Tema 3}{section*.23}{}}
\newlabel{sol:eje:3-no-comprime@cref}{{[subappendix][3][2147483647,1]A.3}{[1][37][]37}}
\newlabel{sol:eje:4-bpb-concepto}{{A.4}{38}{Tema 4}{section*.24}{}}
\newlabel{sol:eje:4-bpb-concepto@cref}{{[subappendix][4][2147483647,1]A.4}{[1][38][]38}}
\newlabel{sol:eje:4-bac}{{A.4}{38}{Tema 4}{section*.25}{}}
\newlabel{sol:eje:4-bac@cref}{{[subappendix][4][2147483647,1]A.4}{[1][38][]38}}
\newlabel{sol:eje:4-descodificar}{{A.4}{39}{Tema 4}{section*.26}{}}
\newlabel{sol:eje:4-descodificar@cref}{{[subappendix][4][2147483647,1]A.4}{[1][39][]39}}
\newlabel{sol:eje:4-abbb}{{A.4}{39}{Tema 4}{section*.27}{}}
\newlabel{sol:eje:4-abbb@cref}{{[subappendix][4][2147483647,1]A.4}{[1][39][]39}}
\newlabel{sol:eje:4-bpb}{{A.4}{40}{Tema 4}{section*.28}{}}
\newlabel{sol:eje:4-bpb@cref}{{[subappendix][4][2147483647,1]A.4}{[1][40][]40}}
\newlabel{sol:eje:4-equivocado}{{A.4}{40}{Tema 4}{section*.29}{}}
\newlabel{sol:eje:4-equivocado@cref}{{[subappendix][4][2147483647,1]A.4}{[1][40][]40}}
\newlabel{sol:eje:4-equivocado-teo}{{A.4}{40}{Tema 4}{section*.30}{}}
\newlabel{sol:eje:4-equivocado-teo@cref}{{[subappendix][4][2147483647,1]A.4}{[1][40][]40}}
\newlabel{sol:eje:5-beta}{{A.5}{41}{Tema 5}{section*.31}{}}
\newlabel{sol:eje:5-beta@cref}{{[subappendix][5][2147483647,1]A.5}{[1][41][]41}}
\newlabel{sol:eje:5-monedas}{{A.5}{41}{Tema 5}{section*.32}{}}
\newlabel{sol:eje:5-monedas@cref}{{[subappendix][5][2147483647,1]A.5}{[1][41][]41}}
\newlabel{sol:eje:5-landauer}{{A.5}{42}{Tema 5}{section*.33}{}}
\newlabel{sol:eje:5-landauer@cref}{{[subappendix][5][2147483647,1]A.5}{[1][42][]42}}
\newlabel{sol:eje:5-rlhf}{{A.5}{42}{Tema 5}{section*.34}{}}
\newlabel{sol:eje:5-rlhf@cref}{{[subappendix][5][2147483647,1]A.5}{[1][42][]42}}
\newlabel{sol:eje:5-lagrange}{{A.5}{42}{Tema 5}{section*.35}{}}
\newlabel{sol:eje:5-lagrange@cref}{{[subappendix][5][2147483647,1]A.5}{[1][42][]42}}
\newlabel{sol:eje:5-dpo}{{A.5}{43}{Tema 5}{section*.36}{}}
\newlabel{sol:eje:5-dpo@cref}{{[subappendix][5][2147483647,1]A.5}{[1][43][]43}}
\makeatother

\begin{document}
\frontmatter
\begin{titlepage}
\centering
\vspace*{2.5cm}
{\large\scshape Anatomía de un LLM}\par
\vspace{0.4cm}
{\normalsize Fundamentos matemáticos de los modelos de lenguaje}\par
\vspace{2.2cm}
\rule{\textwidth}{0.8pt}\par\vspace{0.5cm}
{\small\scshape Fase I \textperiodcentered\ Fundamentos matemáticos}\par\vspace{0.3cm}
{\Huge\bfseries Módulo 6\par}\vspace{0.3cm}
{\LARGE Teoría de la información\par}
\vspace{0.5cm}\rule{\textwidth}{0.4pt}\par
\vspace{1.6cm}
{\large Resumen del módulo}\par
\vfill
{\small Edición de octubre de 2026}\par
\end{titlepage}

\chapter*{Sobre este resumen}
\addcontentsline{toc}{chapter}{Sobre este resumen}
\markboth{Sobre este resumen}{}

Este documento reúne lo indispensable de los apuntes del Módulo 6: el texto expositivo, todas las definiciones y todos los resultados (teoremas, proposiciones, lemas y corolarios), las tablas y las observaciones de contenido. Se han suprimido las demostraciones, los ejemplos, las lecturas, las figuras, los ejercicios y sus soluciones. La numeración de temas, secciones, enunciados y ecuaciones es la misma que en los apuntes, de modo que cualquier referencia puede buscarse allí; las que remiten a elementos suprimidos (un ejemplo, una figura, una ecuación de una demostración) conservan su número en los apuntes.

\section*{Notación}
\begin{center}
\begin{tabular}{@{}ll@{}}
\toprule
Símbolo & Significado \\
\midrule
$\X$, $K=\abs{\X}$ & alfabeto finito y su tamaño \\
$p(x)$, $p(x,y)$, $p(y\mid x)$ & masa de probabilidad, conjunta y condicionada \\
$h(x)=-\log p(x)$ & contenido de información (sorpresa) de $x$ \\
$H(X)=H(p)$ & entropía \\
$H(X,Y)$, $H(Y\mid X)$ & entropía conjunta y condicionada \\
$H(p,q)$ & entropía cruzada \\
$\KL{p}{q}$ & entropía relativa o divergencia de Kullback--Leibler \\
$I(X;Y)$ & información mutua \\
$\ell(x)$, $L=\E[\ell(X)]$ & longitud de la palabra de código de $x$ y longitud media \\
$F(x)$ & probabilidad acumulada de los símbolos anteriores a $x$ \\
$\bpb$ & bits por byte \\
$\pref$, $r$, $\beta$ & política de referencia, recompensa y coeficiente \\
$\log=\log_2$, $\ln$ & logaritmos en base~2 (bits) y natural (nats) \\
\bottomrule
\end{tabular}
\end{center}


\tableofcontents

\mainmatter
% =====================================================================
\chapter{Información y entropía}
\label{tema:1}
% =====================================================================

Este tema define la medida de la información introducida por Shannon~\cite{shannon} y sus versiones para varias variables: entropía, entropía conjunta y entropía condicionada. Se sigue~\cite[§2.1, §2.2 y §2.5]{ct}; la interpretación de la entropía como número de preguntas y la discusión del contenido de información proceden de~\cite[cap.~2 y~4]{mackay} y de la primera clase de~\cite{camit}. En todo el tema, $X$ es una variable aleatoria discreta con valores en un alfabeto finito $\X$ y función de masa $p(x)=P(X=x)$.

\section{El contenido de información de un suceso}

Un suceso seguro no informa de nada; un suceso improbable, al ocurrir, informa mucho. Se busca una función $h$ que asigne a cada probabilidad $p\in(0,1]$ una cantidad de información y que cumpla tres requisitos naturales:
\begin{enumerate}
  \item $h(1)=0$: lo que es seguro no aporta información;
  \item $h$ es decreciente: cuanto menos probable es un suceso, más informa;
  \item $h(pq)=h(p)+h(q)$: la información de dos sucesos independientes es la suma de las informaciones.
\end{enumerate}

\begin{proposicion}[Caracterización del contenido de información]\label{prop:caracterizacion}
Sea $h\colon(0,1]\to\R$ decreciente (en sentido amplio) y tal que $h(pq)=h(p)+h(q)$ para todo $p,q\in(0,1]$. Entonces existe una constante $c\ge 0$ tal que $h(p)=-c\log_2 p$ para todo $p$.
\end{proposicion}

La constante $c$ solo fija la unidad de medida. Se adopta $c=1$.

\begin{definicion}[Contenido de información]\label{def:sorpresa}
El \emph{contenido de información} o \emph{sorpresa} del resultado $x$, con $p(x)>0$, es
\[
h(x)=-\log_2 p(x)\quad\text{bits}.
\]
Con el logaritmo natural la unidad es el \emph{nat}; como $\ln p=\ln 2\cdot\log_2 p$, se tiene $1\text{ nat}=1/\ln 2\approx 1{,}443$ bits.
\end{definicion}

\begin{center}
\begin{tabular}{@{}lcc@{}}
\toprule
Suceso & $p$ & $h$ (bits) \\
\midrule
cara en una moneda equilibrada & $1/2$ & $1$ \\
un~3 en un dado equilibrado & $1/6$ & $2{,}585$ \\
uno entre 8 resultados equiprobables & $1/8$ & $3$ \\
un token al que el modelo asigna $0{,}9$ & $0{,}9$ & $0{,}152$ \\
un token al que el modelo asigna $0{,}1$ & $0{,}1$ & $3{,}322$ \\
uno entre $1024$ resultados equiprobables & $2^{-10}$ & $10$ \\
\bottomrule
\end{tabular}
\end{center}

Si $p(x)=2^{-k}$, el resultado $x$ puede identificarse con $k$ preguntas de respuesta sí o no que dividan por la mitad, en cada paso, el conjunto de candidatos. El contenido de información generaliza esta cuenta a probabilidades arbitrarias~\cite[cap.~4]{mackay}.

\addtocounter{definicion}{1}%

\section{Entropía}

\begin{definicion}[Entropía]\label{def:entropia}
La \emph{entropía} de $X$ es la sorpresa esperada:
\[
H(X)=H(p)=\E\bigl[-\log p(X)\bigr]=-\sum_{x\in\X}p(x)\log p(x),
\]
con el convenio $0\log 0=0$, justificado por $\lim_{t\to 0^{+}}t\log t=0$.
\end{definicion}

La entropía depende solo de las probabilidades, no de los valores que toma $X$. Mide la incertidumbre sobre $X$ antes de observarla o, de forma equivalente, la información que se obtiene, en media, al observarla.

\addtocounter{definicion}{1}%

\addtocounter{definicion}{1}%

\addtocounter{figure}{1}%

\begin{proposicion}[Cotas elementales]\label{prop:cotas-entropia}
Para toda distribución $p$ sobre un alfabeto de $K$ símbolos,
\begin{enumerate}
  \item $H(p)\ge 0$, con igualdad si y solo si $p$ es degenerada, es decir, $p(x_0)=1$ para algún $x_0$;
  \item $H(p)\le\log K$, con igualdad si y solo si $p$ es la distribución uniforme.
\end{enumerate}
\end{proposicion}

\begin{center}
\begin{tabular}{@{}lc@{}}
\toprule
Distribución & Entropía (bits) \\
\midrule
uniforme sobre 27 símbolos (26 letras y el espacio) & $\log 27=4{,}755$ \\
$(1/2,1/4,1/4)$ & $1{,}5$ \\
moneda con $P(\text{cara})=0{,}1$ & $0{,}469$ \\
distribución degenerada & $0$ \\
\bottomrule
\end{tabular}
\end{center}

\section{Entropía conjunta y entropía condicionada}

Sean $X$ e $Y$ variables discretas con función de masa conjunta $p(x,y)$, marginales $p(x)$ y $p(y)$ y condicionada $p(y\mid x)=p(x,y)/p(x)$ cuando $p(x)>0$.

\begin{definicion}[Entropía conjunta y condicionada]\label{def:conjunta}
La \emph{entropía conjunta} de $(X,Y)$ y la \emph{entropía de $Y$ condicionada a $X$} son
\begin{gather*}
H(X,Y)=-\sum_{x,y}p(x,y)\log p(x,y),\\
H(Y\mid X)=\sum_{x}p(x)\,H(Y\mid X=x)=-\sum_{x,y}p(x,y)\log p(y\mid x),
\end{gather*}
donde $H(Y\mid X=x)=-\sum_y p(y\mid x)\log p(y\mid x)$ es la entropía de la distribución condicionada.
\end{definicion}

La entropía condicionada es la incertidumbre que queda sobre $Y$ cuando se conoce $X$, promediada sobre los valores de $X$.

\begin{teorema}[Regla de la cadena para la entropía]\label{teo:cadena-entropia}
$H(X,Y)=H(X)+H(Y\mid X)$.
\end{teorema}

\begin{corolario}[Regla de la cadena para secuencias]\label{cor:cadena-secuencias}
Para variables $X_1,\dots,X_n$,
\[
H(X_1,\dots,X_n)=\sum_{t=1}^{n}H(X_t\mid X_1,\dots,X_{t-1}).
\]
\end{corolario}

El \cref{cor:cadena-secuencias} es la versión informacional de la factorización autorregresiva del módulo~1, $P(x_{1:T})=\prod_t P(x_t\mid x_{1:t-1})$: tomando $-\log$ y esperanzas en esa factorización se obtiene exactamente esta suma. La incertidumbre de un texto completo es la suma de las incertidumbres de cada token dado su contexto.

\addtocounter{definicion}{1}%

\begin{observacion}
En general $H(Y\mid X)\ne H(X\mid Y)$, pero $H(X)-H(X\mid Y)=H(Y)-H(Y\mid X)$, porque ambas diferencias valen $H(X)+H(Y)-H(X,Y)$ por la regla de la cadena aplicada en los dos órdenes. Esta cantidad común es la información mutua del \cref{tema:2}.
\end{observacion}

\section{Tasa de entropía y redundancia del lenguaje}

\begin{definicion}[Tasa de entropía]\label{def:tasa}
La \emph{tasa de entropía} de un proceso $X_1,X_2,\dots$ es
\[
\mathcal{H}=\lim_{n\to\infty}\frac1n H(X_1,\dots,X_n),
\]
cuando el límite existe. Se mide en bits por símbolo.
\end{definicion}

Si las $X_t$ son independientes e idénticamente distribuidas, $H(X_1,\dots,X_n)=nH(X_1)$ y $\mathcal{H}=H(X_1)$. En un texto, en cambio, los símbolos dependen fuertemente del contexto y $\mathcal{H}$ es mucho menor que la entropía de un símbolo aislado~\cite[cap.~4]{ct}. Shannon~\cite{shannon51} estimó la tasa de entropía del inglés pidiendo a personas que adivinaran un texto letra a letra, y obtuvo entre $0{,}6$ y $1{,}3$ bits por letra, muy por debajo de los $\log 27=4{,}755$ bits de un alfabeto de 27 símbolos equiprobables. La diferencia mide la \emph{redundancia} del lenguaje: la parte del texto que es previsible a partir de lo anterior.

\addtocounter{definicion}{1}%


% =====================================================================
\chapter{Entropía relativa e información mutua}
\label{tema:2}
% =====================================================================

Este tema introduce la entropía relativa o divergencia de Kullback--Leibler, que mide el coste de usar una distribución equivocada, y la información mutua, que mide cuánto informa una variable sobre otra. El resultado que lo articula todo es la desigualdad de Gibbs, que se deduce de la desigualdad de Jensen. Se sigue~\cite[§2.3--2.6 y §2.8]{ct}; la interpretación de las dos direcciones de la divergencia en aprendizaje automático procede del tema «Comparing probability distributions and machine learning» de~\cite{camit}, y la de la información mutua, de~\cite[cap.~8]{mackay}.

\section{La desigualdad de Jensen}

\begin{definicion}[Función convexa]\label{def:convexa}
Sea $I\subseteq\R$ un intervalo. Una función $f\colon I\to\R$ es \emph{convexa} si
\[
f\bigl(\lambda a+(1-\lambda)b\bigr)\le\lambda f(a)+(1-\lambda)f(b)\qquad\text{para todo }a,b\in I,\ \lambda\in[0,1],
\]
y \emph{estrictamente convexa} si la desigualdad es estricta cuando $a\ne b$ y $\lambda\in(0,1)$. La función $f$ es (estrictamente) \emph{cóncava} si $-f$ es (estrictamente) convexa.
\end{definicion}

Si $f$ es dos veces derivable con $f''\ge 0$ ($f''>0$) en $I$, entonces es convexa (estrictamente convexa)~\cite[§2.6]{ct}; el módulo~4 estudia este criterio. Así, $\ln$ y $\log$ son estrictamente cóncavas en $(0,\infty)$, porque $(\ln t)''=-1/t^{2}<0$, y $t\mapsto t\log t$ es estrictamente convexa en $(0,\infty)$.

\begin{teorema}[Desigualdad de Jensen]\label{teo:jensen}
Sea $f$ convexa en un intervalo $I$ y sea $X$ una variable aleatoria con valores en $I$ que toma un número finito de valores. Entonces
\[
\E[f(X)]\ge f\bigl(\E[X]\bigr).
\]
Si $f$ es estrictamente convexa, la igualdad se da si y solo si $X$ es constante con probabilidad~1. Para $f$ cóncava la desigualdad se invierte.
\end{teorema}

\section{Entropía relativa}

\begin{definicion}[Entropía relativa]\label{def:kl}
Sean $p$ y $q$ distribuciones sobre $\X$. La \emph{entropía relativa} o \emph{divergencia de Kullback--Leibler} de $p$ respecto de $q$ es
\[
\KL{p}{q}=\sum_{x\in\X}p(x)\log\frac{p(x)}{q(x)}=\E_{X\sim p}\Bigl[\log\frac{p(X)}{q(X)}\Bigr],
\]
con los convenios $0\log\frac{0}{q}=0$ y $p\log\frac{p}{0}=+\infty$ si $p>0$.
\end{definicion}

En particular, $\KL{p}{q}=+\infty$ en cuanto $q$ declara imposible un resultado que, según $p$, puede ocurrir. Es la versión informacional de la pérdida infinita del módulo~1 cuando un modelo asigna probabilidad~0 a un token del texto.

\begin{teorema}[Desigualdad de Gibbs o de la información]\label{teo:gibbs}
Para distribuciones $p$ y $q$ sobre $\X$, $\KL{p}{q}\ge 0$, con igualdad si y solo si $p=q$.
\end{teorema}

\addtocounter{definicion}{1}%

\begin{observacion}
La entropía relativa no es simétrica y no cumple la desigualdad triangular; por eso se llama \emph{divergencia} y no distancia. Sí cumple la propiedad que importa para comparar un modelo con los datos: es no negativa y solo se anula cuando las dos distribuciones coinciden.
\end{observacion}

\section{Entropía cruzada}

\begin{definicion}[Entropía cruzada]\label{def:cruzada}
La \emph{entropía cruzada} de $q$ respecto de $p$ es $H(p,q)=-\sum_{x}p(x)\log q(x)=\E_{X\sim p}[-\log q(X)]$.
\end{definicion}

\begin{proposicion}[Descomposición de la entropía cruzada]\label{prop:descomposicion}
$H(p,q)=H(p)+\KL{p}{q}$. En consecuencia, $H(p,q)\ge H(p)$, con igualdad si y solo si $q=p$.
\end{proposicion}

En el \cref{ej:asimetria}, $H(p)=1{,}5$ bits y $H(p,q)=1{,}822$ bits, cuya diferencia es $\KL{p}{q}=0{,}322$ bits. El \cref{tema:3} dará a esta identidad un sentido operativo: $H(p,q)$ es la longitud media de un código diseñado para $q$ cuando los datos siguen $p$, y $\KL{p}{q}$ son los bits que se desperdician por usar el modelo equivocado.

\addtocounter{definicion}{1}%

\begin{observacion}[Las dos direcciones de la divergencia]\label{obs:direcciones}
Al ajustar una distribución $q$ a una distribución $p$ dada pueden minimizarse dos cantidades distintas. La divergencia \emph{hacia delante} $\KL{p}{q}$ penaliza con fuerza que $q(x)$ sea pequeña donde $p(x)$ no lo es: $q$ tiende a \emph{cubrir} todo el soporte de $p$. Es la que se minimiza al entrenar por máxima verosimilitud. La divergencia \emph{hacia atrás} $\KL{q}{p}$ penaliza que $q$ ponga masa donde $p$ es pequeña: $q$ tiende a \emph{concentrarse} en una de las modas de $p$. Es la que aparece en la inferencia variacional y en la penalización del ajuste por preferencias del \cref{tema:5}~\cite{camit}.
\end{observacion}

\section{Información mutua}

\begin{definicion}[Información mutua]\label{def:mutua}
La \emph{información mutua} entre $X$ e $Y$ es la entropía relativa entre la distribución conjunta y el producto de las marginales:
\[
I(X;Y)=\KL{p(x,y)}{p(x)\,p(y)}=\sum_{x,y}p(x,y)\log\frac{p(x,y)}{p(x)\,p(y)}.
\]
\end{definicion}

\begin{proposicion}\label{prop:mutua-identidades}
\[
I(X;Y)=H(X)+H(Y)-H(X,Y)=H(Y)-H(Y\mid X)=H(X)-H(X\mid Y).
\]
En particular, $I(X;Y)=I(Y;X)$ e $I(X;X)=H(X)$.
\end{proposicion}

La demostración de la primera igualdad se propone en el \cref{eje:2-mutua}; las otras dos se siguen de ella y de la regla de la cadena (\cref{teo:cadena-entropia}). Para $Y=X$ se tiene $H(X\mid X)=0$.

\begin{teorema}\label{teo:mutua-positiva}
$I(X;Y)\ge 0$, con igualdad si y solo si $X$ e $Y$ son independientes. En consecuencia, $H(Y\mid X)\le H(Y)$: condicionar no aumenta, en media, la entropía.
\end{teorema}

\begin{observacion}
La desigualdad $H(Y\mid X)\le H(Y)$ vale en media, no para cada valor: puede ocurrir $H(Y\mid X=x)>H(Y)$ para algún $x$ concreto. Una observación sorprendente puede aumentar la incertidumbre, aunque observar $X$ la reduce en promedio~\cite[§2.6]{ct}.
\end{observacion}

\addtocounter{figure}{1}%

\addtocounter{definicion}{1}%

\addtocounter{definicion}{1}%

\section{La desigualdad del procesamiento de datos}

\begin{definicion}[Información mutua condicionada]
Para variables $X$, $Y$ y $Z$, la \emph{información mutua entre $X$ e $Y$ condicionada a $Z$} es $I(X;Y\mid Z)=H(X\mid Z)-H(X\mid Y,Z)$, donde $H(X\mid Y,Z)$ es la entropía de $X$ condicionada al par $(Y,Z)$.
\end{definicion}

Se cumple $I(X;Y\mid Z)=\sum_z p(z)\,I(X;Y\mid Z=z)\ge 0$, porque cada sumando es la información mutua de la distribución condicionada a $Z=z$ y es no negativo por el \cref{teo:mutua-positiva}. Además, aplicando la regla de la cadena a $H(X)-H(X\mid Y,Z)=\bigl[H(X)-H(X\mid Z)\bigr]+\bigl[H(X\mid Z)-H(X\mid Y,Z)\bigr]$ se obtiene la \emph{regla de la cadena para la información mutua}~\cite[§2.5]{ct}:
\[
I(X;Y,Z)=I(X;Z)+I(X;Y\mid Z).
\]

\begin{definicion}[Cadena de Markov]
Las variables $X$, $Y$ y $Z$ forman una \emph{cadena de Markov} $X\to Y\to Z$ si $p(z\mid x,y)=p(z\mid y)$ para todo $x,y,z$; es decir, si $Z$ depende de $X$ solo a través de $Y$. Esto equivale a que $X$ y $Z$ sean condicionalmente independientes dado $Y$, y en ese caso $I(X;Z\mid Y)=0$.
\end{definicion}

\begin{teorema}[Desigualdad del procesamiento de datos]\label{teo:procesamiento}
Si $X\to Y\to Z$, entonces $I(X;Z)\le I(X;Y)$.
\end{teorema}

\begin{corolario}\label{cor:procesamiento-funcion}
Para toda función $g$, $I(X;g(Y))\le I(X;Y)$.
\end{corolario}

\begin{observacion}
Ningún procesamiento de un dato, por elaborado que sea, aumenta la información que contiene sobre otra variable. Las representaciones internas de una red neuronal son funciones de su entrada; por el \cref{cor:procesamiento-funcion}, no contienen más información sobre la respuesta correcta que la propia entrada. Lo que el entrenamiento consigue es hacer esa información \emph{accesible} para una capa de salida sencilla, no crearla.
\end{observacion}


% =====================================================================
\chapter{Códigos de fuente y código de Huffman}
\label{tema:3}
% =====================================================================

Este tema da a la entropía su significado operativo: es la longitud media mínima, en bits por símbolo, de una descripción sin pérdidas de la fuente. Se estudian los códigos prefijo, la desigualdad de Kraft, las cotas para la longitud óptima y el algoritmo de Huffman~\cite{huffman}. Se sigue~\cite[cap.~5]{ct}, en particular §5.1--5.4, §5.6 y §5.8; MacKay~\cite[cap.~5]{mackay} presenta el mismo material con el nombre de \emph{códigos de símbolos}, que es también el enfoque del tema «Source coding and data compression» de~\cite{camit}.

\section{Códigos de símbolos}

\begin{definicion}[Código]\label{def:codigo}
Un \emph{código binario} para una variable $X$ con alfabeto $\X$ es una aplicación $C\colon\X\to\{\bin{0},\bin{1}\}^{*}$ que asigna a cada símbolo una cadena finita de bits, su \emph{palabra de código}, de longitud $\ell(x)$. La \emph{extensión} de $C$ codifica una secuencia concatenando palabras: $C(x_1\cdots x_n)=C(x_1)\cdots C(x_n)$. La \emph{longitud media} del código es $L(C)=\sum_{x}p(x)\,\ell(x)$.
\end{definicion}

\begin{definicion}[Clases de códigos]\label{def:clases}
Un código es
\begin{enumerate}
  \item \emph{no singular} si $C$ es inyectiva;
  \item \emph{unívocamente descodificable} si su extensión es inyectiva, es decir, si toda cadena de bits procede de, a lo sumo, una secuencia de símbolos;
  \item \emph{prefijo} o \emph{instantáneo} si ninguna palabra de código es prefijo de otra distinta.
\end{enumerate}
\end{definicion}

Cada clase contiene a la siguiente: todo código prefijo es unívocamente descodificable, porque el receptor reconoce el final de cada palabra en cuanto la lee completa, sin mirar los bits posteriores; y todo código unívocamente descodificable es no singular.

\addtocounter{definicion}{1}%

\section{La desigualdad de Kraft}

A cada cadena binaria $c=c_1\cdots c_\ell$ se le asocia el \emph{intervalo diádico} $J(c)=[0.c_1\cdots c_\ell,\ 0.c_1\cdots c_\ell+2^{-\ell})\subseteq[0,1)$, donde $0.c_1\cdots c_\ell=\sum_{i}c_i2^{-i}$. Es el conjunto de números de $[0,1)$ cuya expansión binaria empieza por $c$, y tiene longitud $2^{-\ell}$. Dos intervalos diádicos están contenidos uno en otro o son disjuntos, y $J(c')\subseteq J(c)$ si y solo si $c$ es prefijo de $c'$.

\begin{teorema}[Desigualdad de Kraft]\label{teo:kraft}
Las longitudes $\ell_1,\dots,\ell_m$ de las palabras de un código binario prefijo cumplen
\[
\sum_{i=1}^{m}2^{-\ell_i}\le 1.
\]
Recíprocamente, si unos enteros positivos $\ell_1,\dots,\ell_m$ cumplen esta desigualdad, existe un código prefijo con esas longitudes.
\end{teorema}

\addtocounter{figure}{1}%

\begin{observacion}
McMillan demostró que la desigualdad de Kraft la cumplen también todos los códigos unívocamente descodificables~\cite[§5.5]{ct}. Por tanto, restringirse a códigos prefijo no tiene coste: cualquier conjunto de longitudes alcanzable con un código unívocamente descodificable lo es también con uno prefijo.
\end{observacion}

\section{Cotas para la longitud óptima}

\begin{teorema}[Cota inferior]\label{teo:cota-inferior}
La longitud media de todo código prefijo para $X$ cumple $L\ge H(X)$, con igualdad si y solo si $p(x)=2^{-\ell(x)}$ para todo $x$.
\end{teorema}

La demostración interpreta cada código prefijo como un modelo probabilístico implícito, $r(x)\propto2^{-\ell(x)}$: un código asigna palabras cortas a lo que «cree» probable. Esta lectura es la mitad de la equivalencia entre predecir y comprimir que se completa en el \cref{tema:4}.

\begin{definicion}[Código de Shannon]\label{def:shannon}
Un \emph{código de Shannon} para $p$ es un código prefijo con longitudes $\ell(x)=\lceil-\log p(x)\rceil$.
\end{definicion}

Estas longitudes cumplen la desigualdad de Kraft, porque $2^{-\lceil-\log p(x)\rceil}\le2^{\log p(x)}=p(x)$ y las $p(x)$ suman~1; por el \cref{teo:kraft}, el código existe.

\begin{teorema}[Longitud óptima]\label{teo:cotas}
Sea $L^{*}$ la mínima longitud media de un código prefijo para $X$. Entonces
\[
H(X)\le L^{*}<H(X)+1.
\]
\end{teorema}

El exceso de hasta un bit se debe a que las longitudes han de ser enteras. Se reparte entre muchos símbolos codificando bloques.

\begin{corolario}[Codificación por bloques]\label{cor:bloques}
Sean $X_1,\dots,X_n$ independientes y con la misma distribución que $X$, y sea $L_n^{*}$ la longitud media mínima de un código prefijo para el bloque $(X_1,\dots,X_n)$. Entonces
\[
H(X)\le\frac{L_n^{*}}{n}<H(X)+\frac1n.
\]
En particular, $L_n^{*}/n\to H(X)$ cuando $n\to\infty$.
\end{corolario}

El \cref{cor:bloques} es una forma del \emph{teorema de codificación de fuente} de Shannon: la entropía es el número mínimo de bits por símbolo con que puede describirse la fuente sin pérdidas, y ese mínimo se alcanza asintóticamente~\cite[cap.~4]{mackay}.

\section{El código de Huffman}

\begin{definicion}[Algoritmo de Huffman]\label{def:huffman}
Dada una distribución sobre $K\ge 2$ símbolos, el \emph{algoritmo de Huffman} construye un árbol binario: mientras quede más de un nodo, fusiona los dos nodos de menor probabilidad en un nodo nuevo cuya probabilidad es la suma de ambas. Al terminar, las ramas que salen de cada nodo se etiquetan con \bin{0} y \bin{1}, y la palabra de cada símbolo es la secuencia de etiquetas desde la raíz hasta su hoja.
\end{definicion}

Cada fusión añade un bit a la palabra de todos los símbolos de los dos nodos fusionados; por eso los símbolos poco probables, que se fusionan pronto, reciben palabras largas.

\addtocounter{definicion}{1}%

\addtocounter{figure}{1}%

La optimalidad del algoritmo se apoya en dos propiedades de los códigos óptimos.

\begin{lema}\label{lem:optimo}
Sea $p_1\ge p_2\ge\dots\ge p_K>0$, con $K\ge 2$. Existe un código prefijo óptimo con longitudes $\ell_1,\dots,\ell_K$ tal que
\begin{enumerate}
  \item $\ell_1\le\ell_2\le\dots\le\ell_K$;
  \item las palabras de los dos símbolos menos probables, $K-1$ y $K$, tienen la longitud máxima $\ell_K$ y solo difieren en el último bit.
\end{enumerate}
\end{lema}

\begin{teorema}[Optimalidad del código de Huffman]\label{teo:huffman}
El código de Huffman tiene longitud media mínima entre todos los códigos prefijo para $p$.
\end{teorema}

Por los Teoremas~\ref{teo:cotas} y~\ref{teo:huffman}, $H(p)\le L_{\text{Huffman}}<H(p)+1$, con $L_{\text{Huffman}}=H(p)$ cuando todas las probabilidades son potencias de $1/2$. El código óptimo no es único: con otros desempates en el \cref{ej:huffman} se obtiene un árbol distinto (\cref{eje:3-dos-huffman}).

\section{El límite de los códigos de símbolos}

Un código de símbolos gasta un número entero de bits por símbolo, y al menos un bit. Si $p=(0{,}9;\ 0{,}1)$, la entropía es $H_b(0{,}1)=0{,}469$ bits, pero cualquier código de símbolos para dos resultados gasta 1 bit por símbolo: más del doble. La codificación por bloques del \cref{cor:bloques} reduce el exceso, pero un bloque de $n$ símbolos sobre un alfabeto de $K$ símbolos tiene $K^{n}$ valores posibles, y la tabla del código crece exponencialmente.

\addtocounter{definicion}{1}%


% =====================================================================
\chapter{Codificación aritmética: predecir es comprimir}
\label{tema:4}
% =====================================================================

Este tema presenta la codificación aritmética, que representa un mensaje completo mediante un subintervalo de $[0,1)$ de anchura igual a su probabilidad, y demuestra que con ella cualquier modelo probabilístico se convierte en un compresor cuya longitud es, salvo dos bits, la log-verosimilitud negativa del mensaje en bits. La construcción sigue la codificación de Shannon--Fano--Elias de~\cite[§5.9]{ct} y su extensión a secuencias de~\cite[§13.3]{ct}; la presentación como \emph{stream code} con un modelo adaptativo es la de MacKay~\cite[cap.~6]{mackay}, y la implementación práctica, la de Witten, Neal y Cleary~\cite{wnc}. La conexión con los modelos de lenguaje sigue a Delétang \emph{et al.}~\cite{deletang} y a Gao \emph{et al.}~\cite{pile}.

\section{Intervalos y probabilidades}

Sea $\X=\{x^{(1)},\dots,x^{(K)}\}$ un alfabeto con un orden fijado y sea $q$ un modelo autorregresivo, que asigna a cada contexto $x_{1:t-1}$ una distribución $q(\cdot\mid x_{1:t-1})$ sobre $\X$. La \emph{probabilidad acumulada} de $x$ en la posición $t$ es la masa de los símbolos anteriores en el orden fijado:
\[
F_t(x)=\sum_{x'<x}q(x'\mid x_{1:t-1}).
\]

\begin{definicion}[Intervalo de una secuencia]\label{def:intervalo}
El \emph{intervalo} de la secuencia $x_{1:n}$ es $I(x_{1:n})=[a_n,a_n+w_n)$, donde $a_0=0$, $w_0=1$ y, para $t=1,\dots,n$,
\[
a_t=a_{t-1}+w_{t-1}\,F_t(x_t),\qquad w_t=w_{t-1}\,q(x_t\mid x_{1:t-1}).
\]
\end{definicion}

Cada paso divide el intervalo actual en $K$ subintervalos consecutivos, de anchuras proporcionales a las probabilidades del símbolo siguiente, y se queda con el del símbolo observado.

\begin{proposicion}\label{prop:intervalos}
Para todo $n\ge 0$:
\begin{enumerate}
  \item la anchura de $I(x_{1:n})$ es $w_n=q(x_{1:n})=\prod_{t=1}^{n}q(x_t\mid x_{1:t-1})$;
  \item los intervalos de todas las secuencias de longitud $n$ forman una partición de $[0,1)$;
  \item $I(x_{1:n})\subseteq I(x_{1:n-1})$ si $n\ge 1$.
\end{enumerate}
\end{proposicion}

\addtocounter{definicion}{1}%

\addtocounter{figure}{1}%

\section{Palabras de código y longitud}

Para transmitir el intervalo basta enviar una cadena binaria $c$ cuyo intervalo diádico $J(c)$ (\cref{tema:3}) esté contenido en $I(x_{1:n})$: el receptor sabrá entonces que el número enviado, cualesquiera que sean los bits que lo sigan, está dentro del intervalo del mensaje.

\begin{teorema}\label{teo:palabra}
Sea $I=[a,a+w)\subseteq[0,1)$ con $w>0$, y sea $\ell=\lceil-\log w\rceil+1$. Existe una cadena binaria $c$ de longitud $\ell$ con $J(c)\subseteq I$.
\end{teorema}

\begin{corolario}[Longitud de la codificación aritmética]\label{cor:longitud}
Para cada $n$, asignar a cada secuencia $x_{1:n}$ con $q(x_{1:n})>0$ la cadena del \cref{teo:palabra} para $I(x_{1:n})$ define un código prefijo sobre las secuencias de longitud $n$, con
\[
\ell(x_{1:n})=\lceil-\log q(x_{1:n})\rceil+1<-\log q(x_{1:n})+2.
\]
\end{corolario}

\addtocounter{definicion}{1}%

\paragraph{Descodificación.} El receptor conoce el modelo $q$ y la longitud $n$ (o bien el alfabeto incluye un símbolo de fin, como el símbolo \fin\ del módulo~1). Recibida $c$, para $t=1,\dots,n$ subdivide el intervalo actual según $q(\cdot\mid x_{1:t-1})$, determina el único símbolo $x_t$ cuyo subintervalo contiene a $J(c)$ y lo añade al contexto. Por la \cref{prop:intervalos}, ese símbolo existe, es único y coincide con el enviado. En la práctica, la aritmética de precisión infinita se sustituye por aritmética entera con renormalización, que emite los bits a medida que se fijan, con un sobrecoste despreciable~\cite{wnc}.

\section{Predecir es comprimir}

\begin{teorema}[Compresión con un modelo]\label{teo:compresion}
Sea $q$ un modelo autorregresivo cualquiera. La codificación aritmética con $q$ codifica cada secuencia $x_{1:n}$ con
\[
\ell(x_{1:n})<-\log q(x_{1:n})+2=\sum_{t=1}^{n}-\log q(x_t\mid x_{1:t-1})+2\quad\text{bits}.
\]
Si $X_{1:n}$ sigue una distribución $p$, la longitud media cumple
\[
H(X_{1:n})\le\E[\ell(X_{1:n})]<H(p,q)+2=H(X_{1:n})+\KL{p}{q}+2,
\]
donde $p$ y $q$ denotan aquí las distribuciones de la secuencia completa.
\end{teorema}

\begin{corolario}\label{cor:por-simbolo}
Si las $X_t$ son independientes con distribución $p$ y el modelo es sin memoria con distribución $q$, la longitud media por símbolo cumple
\[
H(p)\le\frac{\E[\ell(X_{1:n})]}{n}<H(p)+\KL{p}{q}+\frac{2}{n}.
\]
Con $q=p$, la codificación aritmética alcanza la entropía con un sobrecoste de $2/n$ bits por símbolo, sin tablas que crezcan con $n$.
\end{corolario}

El \cref{teo:compresion} es la mitad de la equivalencia anunciada: \emph{todo predictor da un compresor}, y el tamaño del archivo comprimido es la suma de las sorpresas del modelo. La otra mitad se vio en el \cref{teo:cota-inferior}: \emph{todo compresor da un predictor}, porque las longitudes de un código prefijo definen la distribución $r(x)\propto2^{-\ell(x)}$. Mejorar uno equivale a mejorar el otro.

\addtocounter{figure}{1}%

\addtocounter{definicion}{1}%

\section{Bits por byte}

La pérdida por token y la perplejidad dependen del tokenizador: un vocabulario mayor produce menos tokens, cada uno más difícil de predecir, y la cifra por token cambia aunque el modelo comprima igual de bien. Para comparar modelos con tokenizadores distintos se normaliza por la longitud del texto, que no depende del tokenizador.

\begin{definicion}[Bits por byte]\label{def:bpb}
Sea un texto de $B$ bytes en la codificación UTF-8, dividido por el tokenizador en $N$ tokens. Los \emph{bits por byte} del modelo $q$ sobre el texto son
\[
\bpb=\frac{1}{B}\sum_{t=1}^{N}-\log_2 q(x_t\mid x_{1:t-1}).
\]
\end{definicion}

\begin{proposicion}\label{prop:bpb}
Si $\mathcal{J}$ es la pérdida media en nats por token y $\bar b=B/N$ el número medio de bytes por token, entonces $\bpb=\mathcal{J}/(\bar b\ln 2)$.
\end{proposicion}

Por el \cref{teo:compresion}, el bpb es, salvo un término despreciable, el tamaño del texto comprimido con el modelo dividido por su tamaño original en bytes; un bpb de $b$ significa que el texto se comprime a una fracción $b/8$ de su tamaño.

\addtocounter{definicion}{1}%

\begin{center}
\begin{tabular}{@{}llcl@{}}
\toprule
Métrica & Unidad & Tokenizador & Uso habitual \\
\midrule
pérdida & nats por token & depende & entrenamiento \\
perplejidad & $e^{\text{pérdida}}$ & depende & número efectivo de opciones \\
bits por carácter & bits por carácter & no depende & modelos de caracteres \\
bits por byte & bits por byte & no depende & comparar modelos distintos \\
\bottomrule
\end{tabular}
\end{center}

Los bits por carácter tampoco dependen del tokenizador, pero sí de cómo se cuentan los caracteres; los bytes UTF-8 están definidos sin ambigüedad. En español, las vocales con tilde, la \emph{ñ} y las comillas angulares ocupan 2 bytes cada una, de modo que un mismo texto tiene más bytes que caracteres y su bpb es algo menor que sus bits por carácter. El artículo del conjunto de datos The Pile~\cite{pile} popularizó el bpb para comparar modelos con tokenizadores distintos. Anthropic, Google y xAI no publican el bpb de sus modelos actuales.


% =====================================================================
\chapter{Entropía, física y ajuste por preferencias}
\label{tema:5}
% =====================================================================

Este tema relaciona la entropía de Shannon con la entropía de la física estadística y muestra que la misma herramienta, la desigualdad de Gibbs, resuelve tres problemas en apariencia distintos: la distribución de máxima entropía con una restricción de energía media, que es la distribución de Boltzmann y, en particular, la función softmax; el principio variacional de la energía libre; y el problema de maximizar una recompensa con una penalización de Kullback--Leibler, que está en la base del ajuste de los LLM por preferencias humanas. Se sigue~\cite[cap.~12]{ct} y el artículo de Jaynes~\cite{jaynes} para la máxima entropía, el artículo de Landauer~\cite{landauer} para el coste físico del borrado y los artículos de Ouyang \emph{et al.}~\cite{instructgpt} y Rafailov \emph{et al.}~\cite{dpo} para el ajuste por preferencias. El tema «Applications of information theory in other sciences» de~\cite{camit} trata el mismo puente con la física.

En este tema los logaritmos son naturales y las entropías y divergencias se miden en nats, salvo indicación expresa.

\section{La entropía en la física estadística}

Un sistema físico macroscópico, descrito por unas pocas magnitudes (energía, volumen, número de partículas), es compatible con un número enorme de configuraciones microscópicas o \emph{microestados}.

\begin{definicion}[Entropías de Boltzmann y de Gibbs]\label{def:boltzmann}
Si un macroestado es compatible con $W$ microestados equiprobables, su \emph{entropía de Boltzmann} es $S=k_B\ln W$. Si los microestados tienen probabilidades $p_i$, la \emph{entropía de Gibbs} es
\[
S=-k_B\sum_ip_i\ln p_i,
\]
donde $k_B=1{,}380649\cdot10^{-23}\,\mathrm{J/K}$ es la constante de Boltzmann.
\end{definicion}

La entropía de Gibbs es $k_B$ veces la entropía de Shannon en nats, y la de Boltzmann es el caso uniforme, en el que $H=\ln W$. La entropía termodinámica de un macroestado es, salvo la constante $k_B$, la cantidad de información que falta para identificar el microestado~\cite{jaynes}.

\begin{center}
\begin{tabular}{@{}ll@{}}
\toprule
Física estadística & Teoría de la información \\
\midrule
microestado & mensaje o secuencia concreta \\
$S=k_B\ln W$ & $H=\log W$ bits (caso uniforme) \\
$S=-k_B\sum_ip_i\ln p_i$ & $H=-\sum_ip_i\log p_i$ \\
distribución de Boltzmann $e^{-E/k_BT}/Z$ & softmax de los logits con temperatura \\
energía libre $F=\langle E\rangle-TS$ & objetivo con penalización de Kullback--Leibler \\
\bottomrule
\end{tabular}
\end{center}

\begin{proposicion}\label{prop:binomial}
Para $0\le k\le n$, $\binom{n}{k}\le 2^{nH_b(k/n)}$, donde $H_b$ es la entropía binaria en bits.
\end{proposicion}

\addtocounter{definicion}{1}%

\section{Máxima entropía y distribución de Boltzmann}

Cuando de una distribución solo se conoce el valor medio de una magnitud, el principio de máxima entropía de Jaynes~\cite{jaynes} propone elegir, entre todas las distribuciones compatibles, la de mayor entropía: la que no añade ninguna información que no se tenga.

\begin{teorema}[Distribución de máxima entropía]\label{teo:maxent}
Sea $E\colon\X\to\R$ una función de energía sobre un alfabeto finito, sea $\lambda\in\R$ y sea
\[
p_\lambda(x)=\frac{e^{-\lambda E(x)}}{Z(\lambda)},\qquad Z(\lambda)=\sum_{x\in\X}e^{-\lambda E(x)}.
\]
Para toda distribución $p$ sobre $\X$ con la misma energía media, $\E_p[E]=\E_{p_\lambda}[E]$, se cumple $H(p)\le H(p_\lambda)$, con igualdad si y solo si $p=p_\lambda$.
\end{teorema}

En física, $\lambda=1/(k_BT)$ y $p_\lambda$ es la \emph{distribución de Boltzmann} a temperatura $T$: los estados de energía baja son más probables, y tanto más cuanto menor es la temperatura.

\addtocounter{definicion}{1}%

\section{Energía libre y principio variacional}

El problema siguiente aparece en física, con la energía libre, y en el ajuste de modelos de lenguaje, con una recompensa. Sean $\pref$ una distribución de referencia sobre un conjunto finito de respuestas $\mathcal{Y}$, con $\pref(y)>0$ para todo $y$, una recompensa $r\colon\mathcal{Y}\to\R$ y un coeficiente $\beta>0$. Se busca la distribución $\pi$ que maximiza
\[
\mathcal{F}(\pi)=\E_{y\sim\pi}[r(y)]-\beta\,\KL{\pi}{\pref}.
\]

\begin{teorema}[Principio variacional de Gibbs]\label{teo:variacional}
Sean $Z=\sum_y\pref(y)\,e^{r(y)/\beta}$ y $\pi^{*}(y)=\pref(y)\,e^{r(y)/\beta}/Z$. Para toda distribución $\pi$ sobre $\mathcal{Y}$,
\[
\mathcal{F}(\pi)=\beta\ln Z-\beta\,\KL{\pi}{\pi^{*}}.
\]
En consecuencia, $\pi^{*}$ es el único maximizador de $\mathcal{F}$ y el valor máximo es $\beta\ln Z$.
\end{teorema}

\begin{corolario}\label{cor:energia-libre}
Si $\pref$ es uniforme sobre $K$ respuestas, $\mathcal{F}(\pi)=\E_\pi[r]+\beta H(\pi)-\beta\ln K$. Con energía $E=-r$ y temperatura $\beta$, maximizar $\mathcal{F}$ equivale a minimizar la \emph{energía libre} $\E_\pi[E]-\beta H(\pi)$, cuyo minimizador es la distribución de Boltzmann.
\end{corolario}

\section{El ajuste por preferencias con penalización de Kullback--Leibler}

Tras el preentrenamiento y el ajuste supervisado, un modelo de lenguaje se ajusta a las preferencias humanas con aprendizaje por refuerzo (RLHF, módulo~21). Un \emph{modelo de recompensa} $r(x,y)$, entrenado con comparaciones humanas, puntúa la respuesta $y$ al mensaje $x$. Maximizar solo $r$ llevaría el modelo hacia respuestas extrañas que explotan los defectos del modelo de recompensa (\emph{reward hacking}). Por ello, Ouyang \emph{et al.}~\cite{instructgpt} añaden a la recompensa una penalización por alejarse del modelo supervisado $\pref$, que en esperanza es $\beta\,\KL{\pi}{\pref}$; para cada mensaje $x$ se maximiza la cantidad $\mathcal{F}$ del apartado anterior.

Por el \cref{teo:variacional}, la política óptima es $\pi^{*}(y\mid x)\propto\pref(y\mid x)\,e^{r(x,y)/\beta}$: la referencia «recalentada» por la recompensa, con $\beta$ en el papel de temperatura. En la práctica el conjunto de respuestas es inmenso y $\pi^{*}$ no se calcula de forma explícita, sino que se aproxima con métodos de gradiente; el teorema describe a qué apunta esa optimización.

\addtocounter{definicion}{1}%

\begin{observacion}[Optimización directa de preferencias]
Despejando la recompensa en la solución óptima, $r(x,y)=\beta\ln\bigl(\pi^{*}(y\mid x)/\pref(y\mid x)\bigr)+\beta\ln Z(x)$: cada política define implícitamente una recompensa. Rafailov \emph{et al.}~\cite{dpo} observan que, en el modelo de Bradley--Terry de comparaciones por pares, el término $\beta\ln Z(x)$ se cancela, y obtienen un método, DPO, que ajusta la política directamente con las comparaciones humanas, sin entrenar un modelo de recompensa separado (\cref{eje:5-dpo}).
\end{observacion}

\section{El coste físico de la información}

\begin{proposicion}[Principio de Landauer]\label{prop:landauer}
Borrar un bit de información en un sistema en contacto con un entorno a temperatura $T$ disipa en el entorno un calor de al menos $k_BT\ln 2$.
\end{proposicion}

La idea, debida a Landauer~\cite{landauer}, es la siguiente. Un bit de memoria cuyo valor se desconoce está en uno de dos estados equiprobables, con entropía $k_B\ln 2$; tras el borrado, que lo pone siempre a~0, su entropía es nula. Como la entropía total no puede disminuir, al menos $k_B\ln 2$ de entropía pasa al entorno, lo que exige un calor $Q\ge T\,\Delta S=k_BT\ln 2$. A $T=300$~K, este mínimo vale $1{,}380649\cdot10^{-23}\cdot300\cdot0{,}6931=2{,}87\cdot10^{-21}$~J por bit. Los circuitos actuales disipan muchos órdenes de magnitud más por operación; el módulo~17 retoma la cuestión del coste energético del cómputo.



\backmatter
% =====================================================================
\begin{thebibliography}{99}
\addcontentsline{toc}{chapter}{Bibliografía}
\markboth{Bibliografía}{}
\small

\bibitem{ct}
T.~M. Cover y J.~A. Thomas.
\newblock \emph{Elements of Information Theory}, 2.ª ed.
\newblock Wiley, 2006.

\bibitem{deletang}
G.~Delétang, A.~Ruoss, P.-A. Duquenne, E.~Catt, T.~Genewein, C.~Mattern, J.~Grau-Moya, L.~K. Wenliang, M.~Aitchison, L.~Orseau, M.~Hutter y J.~Veness.
\newblock Language modeling is compression.
\newblock \emph{arXiv:2309.10668}, 2023.
\newblock \url{https://arxiv.org/abs/2309.10668}

\bibitem{pile}
L.~Gao, S.~Biderman, S.~Black \emph{et al.}
\newblock The Pile: An 800GB dataset of diverse text for language modeling.
\newblock \emph{arXiv:2101.00027}, 2020.
\newblock \url{https://arxiv.org/abs/2101.00027}

\bibitem{camit}
R.~Harle.
\newblock Information Theory.
\newblock University of Cambridge, Department of Computer Science and Technology, Part~II, curso 2025--26.
\newblock \emph{Examples Sheet}: \url{https://www.cl.cam.ac.uk/teaching/2526/InfoTheory/InfoTheoryExamplesSheet.pdf}

\bibitem{chinchilla}
J.~Hoffmann, S.~Borgeaud, A.~Mensch \emph{et al.}
\newblock Training compute-optimal large language models.
\newblock \emph{arXiv:2203.15556}, 2022.
\newblock \url{https://arxiv.org/abs/2203.15556}

\bibitem{huffman}
D.~A. Huffman.
\newblock A method for the construction of minimum-redundancy codes.
\newblock \emph{Proceedings of the IRE}, 40(9):1098--1101, 1952.

\bibitem{jaynes}
E.~T. Jaynes.
\newblock Information theory and statistical mechanics.
\newblock \emph{Physical Review}, 106(4):620--630, 1957.

\bibitem{landauer}
R.~Landauer.
\newblock Irreversibility and heat generation in the computing process.
\newblock \emph{IBM Journal of Research and Development}, 5(3):183--191, 1961.

\bibitem{mackay}
D.~J.~C. MacKay.
\newblock \emph{Information Theory, Inference, and Learning Algorithms}.
\newblock Cambridge University Press, 2003.
\newblock \url{http://www.inference.org.uk/itila/book.html}

\bibitem{instructgpt}
L.~Ouyang, J.~Wu, X.~Jiang \emph{et al.}
\newblock Training language models to follow instructions with human feedback.
\newblock \emph{arXiv:2203.02155}, 2022.
\newblock \url{https://arxiv.org/abs/2203.02155}

\bibitem{mit6441}
Y.~Polyanskiy.
\newblock 6.441: Information Theory.
\newblock Massachusetts Institute of Technology, MIT OpenCourseWare, primavera de 2016.
\newblock \url{https://ocw.mit.edu/courses/6-441-information-theory-spring-2016/}

\bibitem{dpo}
R.~Rafailov, A.~Sharma, E.~Mitchell, S.~Ermon, C.~D. Manning y C.~Finn.
\newblock Direct preference optimization: Your language model is secretly a reward model.
\newblock \emph{arXiv:2305.18290}, 2023.
\newblock \url{https://arxiv.org/abs/2305.18290}

\bibitem{shannon}
C.~E. Shannon.
\newblock A mathematical theory of communication.
\newblock \emph{Bell System Technical Journal}, 27:379--423 y 623--656, 1948.

\bibitem{shannon51}
C.~E. Shannon.
\newblock Prediction and entropy of printed English.
\newblock \emph{Bell System Technical Journal}, 30(1):50--64, 1951.

\bibitem{wnc}
I.~H. Witten, R.~M. Neal y J.~G. Cleary.
\newblock Arithmetic coding for data compression.
\newblock \emph{Communications of the ACM}, 30(6):520--540, 1987.

\end{thebibliography}


\end{document}
